\documentclass[11pt,a4paper]{article}

\usepackage[margin=1in]{geometry}
\usepackage{amsmath,amssymb,amsfonts,bm}
\usepackage{booktabs}
\usepackage[numbers,sort&compress]{natbib}
\usepackage{hyperref}
\usepackage{microtype}
\usepackage{enumitem}

\usepackage{url}
\hypersetup{
  colorlinks=true,
  linkcolor=blue,
  urlcolor=blue,
  citecolor=blue
}
\newcommand{\doiref}[1]{\href{https://doi.org/#1}{doi:\nolinkurl{#1}}}

\title{Quasistatic Sphere and Molecular Point Dipole:\\
\large An Analytical Model of Hybrid Absorption}
\author{}
\date{\today}

\begin{document}
\maketitle

\begin{abstract}
\noindent
A classical analytical model is developed for the linear optical response of a
spherical metal nanoparticle coupled to a single molecular (or atomic)
transition treated as a point polarizable dipole.  The nanoparticle is described
in the quasistatic Mie limit by a frequency-dependent dielectric function
$\varepsilon(\omega)$.  Local-field enhancement and multipolar image response
of the sphere yield a closed-form effective molecular polarizability
$\alpha_{\mathrm{eff}}(\omega)$ for radial (parallel) and tangential
(perpendicular) orientations~\citep{gersten1980,gersten1981,bohren1983}.
Absorption is obtained from $-\omega\,\operatorname{Im}\alpha_{\mathrm{eff}}(\omega)$.
A concrete parameterization uses the Raki\'{c} Drude--Lorentz representation of
gold~\citep{rakic1998} as implemented in the Meep material \texttt{Au}~\citep{oskooi2010},
a non-dispersive aqueous host with refractive index $n=1.33$ ($\varepsilon_b=n^2$),
and a single Lorentz oscillator for the sodium $D$-line transition.
\end{abstract}

\tableofcontents
\newpage

% =============================================================================
\section{Physical system}
\label{sec:system}
% =============================================================================

Consider a sphere of radius $a$ centered at the origin, with isotropic dielectric
function $\varepsilon(\omega)$, embedded in a homogeneous host medium of real,
frequency-independent permittivity $\varepsilon_b$ (relative to vacuum).  A
molecule (or atom) is modeled as a point dipole of polarizability
$\alpha_m(\omega)$ located at distance $d>a$ from the center, so that the
surface gap is $d-a$.  For aqueous surroundings with refractive index
$n=1.33$ one has $\varepsilon_b=n^2\approx 1.769$; vacuum corresponds to
$\varepsilon_b=1$.\\

\noindent
Two principal orientations are distinguished relative to the sphere--molecule
axis $\hat{\mathbf{r}}$:
\begin{itemize}[leftmargin=1.4em]
  \item \textbf{Parallel (radial):} the molecular dipole and the driving field
        lie along $\hat{\mathbf{r}}$.
  \item \textbf{Perpendicular (tangential):} both lie in a plane orthogonal to
        $\hat{\mathbf{r}}$.
\end{itemize}

\noindent
The quasistatic approximation is adopted throughout: the sphere diameter is much
smaller than the optical wavelength, retardation across the particle is neglected,
and the incident field is taken as spatially uniform over the sphere~\citep{bohren1983,novotny2012}.
For gold spheres of radius tens of nanometers in the visible, this is a controlled first
approximation; radiative damping and retardation shift the true localized surface
plasmon resonance (LSPR) relative to the quasistatic Fr\"{o}hlich condition derived
below~\citep{maier2007,bohren1983}.  All multipole and local-field formulas below are written for general
$\varepsilon_b$ and are specialized to water by setting $n=1.33$.

% =============================================================================
\section{Dielectric response of the sphere}
\label{sec:sphere}
% =============================================================================

\subsection{Quasistatic multipole coefficients}

In the quasistatic limit of Mie theory~\citep{mie1908,bohren1983}, the response of a sphere in a host of permittivity
$\varepsilon_b$ to an external potential expandable in spherical harmonics of
order $\ell$ is fully characterized by the multipole polarizabilities (volume form)
\begin{equation}
  \beta_\ell(\omega)
  =
  a^{2\ell+1}\,
  \frac{\varepsilon(\omega)-\varepsilon_b}{\varepsilon(\omega)+(\ell+1)\varepsilon_b/\ell},
  \qquad
  \ell = 1,2,\ldots
  \label{eq:beta}
\end{equation}
The electric dipole channel is
\begin{equation}
  \beta_1(\omega)
  =
  a^3\,
  \frac{\varepsilon(\omega)-\varepsilon_b}{\varepsilon(\omega)+2\varepsilon_b}.
  \label{eq:beta1}
\end{equation}
A uniform external field in the medium therefore induces a sphere dipole
proportional to $\beta_1\mathbf{E}_{\mathrm{inc}}$ (in the present units).
Vacuum formulas are recovered by setting $\varepsilon_b=1$.

\subsection{Fr\"{o}hlich resonance}

The dipole denominator vanishes (in the absence of losses) when
\begin{equation}
  \operatorname{Re}\varepsilon(\omega_{\mathrm{LSPR}}) = -2\,\varepsilon_b,
  \label{eq:frohlich}
\end{equation}
which is the Fr\"{o}hlich condition for a sphere in a dielectric host~\citep{frohlich1958,bohren1983,maier2007}.  For
water with $n=1.33$,
\begin{equation}
  \varepsilon_b = n^2 \approx 1.769
  \qquad\Rightarrow\qquad
  \operatorname{Re}\varepsilon(\omega_{\mathrm{LSPR}}) \approx -3.54,
\end{equation}
red-shifting the quasistatic LSPR relative to the vacuum condition
$\operatorname{Re}\varepsilon=-2$.  With realistic $\operatorname{Im}\varepsilon>0$
the resonance is a peak in $|\beta_1(\omega)|$ near that energy rather than a true
pole on the real axis.

% =============================================================================
\section{Local field and reflected field at the molecule}
\label{sec:fields}
% =============================================================================

\subsection{Incident-field enhancement}

A uniform incident field $\mathbf{E}_{\mathrm{inc}}$ (the field far from the
sphere in the host medium) polarizes the sphere.  On the radial axis at distance
$d$, the total field \emph{in the absence of the molecule} is enhanced by the
dipole term~\citep{bohren1983,novotny2012}
\begin{align}
  G_\parallel(\omega)
  &=
  1 + 2
  \left(\frac{a}{d}\right)^3
  \frac{\varepsilon-\varepsilon_b}{\varepsilon+2\varepsilon_b}
  =
  1 + 2\,\frac{\beta_1}{d^3},
  \label{eq:Gpar} \\[0.35em]
  G_\perp(\omega)
  &=
  1 -
  \left(\frac{a}{d}\right)^3
  \frac{\varepsilon-\varepsilon_b}{\varepsilon+2\varepsilon_b}
  =
  1 - \frac{\beta_1}{d^3}.
  \label{eq:Gperp}
\end{align}
Near the LSPR, $|G_\parallel|$ can substantially exceed unity when $a/d$ is not
small.

\subsection{Image multipoles and self-interaction}

A molecular dipole $\mathbf{p}$ induces higher multipoles in the sphere, which
radiate a reflected field $\mathbf{E}_{\mathrm{refl}}=S\,\mathbf{p}$ back at the
molecular site.  With polarizability volumes, the radial and tangential
propagators admit the multipole series~\citep{gersten1980,gersten1981,ruppin1982}
\begin{align}
  S_\parallel(\omega)
  &=
  \sum_{\ell=1}^{\ell_{\max}}
  (\ell+1)^2\,
  \frac{\beta_\ell(\omega)}{d^{2\ell+4}},
  \label{eq:Spar} \\[0.35em]
  S_\perp(\omega)
  &=
  \sum_{\ell=1}^{\ell_{\max}}
  \frac{\ell(\ell+1)}{2}\,
  \frac{\beta_\ell(\omega)}{d^{2\ell+4}}.
  \label{eq:Sperp}
\end{align}
Equivalently, in terms of the ratio $a/d<1$,
\begin{equation}
  \frac{\beta_\ell}{d^{2\ell+4}}
  =
  \frac{1}{d^3}
  \left(\frac{a}{d}\right)^{2\ell+1}
  \frac{\varepsilon-\varepsilon_b}{\varepsilon+(\ell+1)\varepsilon_b/\ell},
\end{equation}
so that the series is built from successive powers of $(a/d)^2$ and remains
numerically stable even for large $\ell_{\max}$.  When the gap $d-a$ is small,
many multipoles contribute and $\ell_{\max}$ must be increased until $S$ converges.

% =============================================================================
\section{Molecular polarizability and effective response}
\label{sec:molecule}
% =============================================================================

\subsection{Bare molecule}

The molecule is assigned an isotropic Lorentzian polarizability volume
\begin{equation}
  \alpha_m(\omega)
  =
  \alpha_0\,
  \frac{\omega_m^2}{\omega_m^2 - \omega^2 - i\gamma_m\omega},
  \label{eq:alpham}
\end{equation}
where $\omega_m$ is the transition energy, $\gamma_m$ a phenomenological
linewidth, and $\alpha_0$ the static polarizability volume
$\alpha_m(0)=\alpha_0$.  This is the classical point-dipole analogue of a
two-level resonance~\citep{jackson1999,novotny2012}.

\subsection{Dressed (effective) polarizability}
\label{sec:alphaeff-deriv}

We now derive the effective molecular polarizability
\begin{equation}
  \alpha_{\mathrm{eff}}(\omega)
  \equiv
  \frac{p}{E_{\mathrm{inc}}}
  =
  \frac{\alpha_m(\omega)\, G(\omega)}{1 - \alpha_m(\omega)\, S(\omega)}
  \label{eq:alphaeff}
\end{equation}
from the classical coupled-dipole picture of Gersten and Nitzan~\citep{gersten1980},
and then identify $G$ and $S$ for a multipolar sphere~\citep{gersten1981,bohren1983}.

\paragraph{Step 1: two polarizable bodies in an external field.}
Following Ref.~\citenum{gersten1980}, consider two classical objects in a uniform
incident field $\mathbf{E}_{\mathrm{inc}}\equiv\mathbf{E}_0$ (the quasistatic part of the
incident radiation).  The molecule (label 1) has polarizability $\boldsymbol{\alpha}_1=\alpha_m$
(taken isotropic and scalar along the chosen axis for simplicity).  The second body
(label 2)---later the nanoparticle---has polarizability $\boldsymbol{\alpha}_2$.
Each dipole responds to the total electric field at its location:
\begin{align}
  \boldsymbol{\mu}_1
  &=
  \boldsymbol{\alpha}_1
  \bigl(
    \mathbf{E}_0
    +
    \mathbf{M}(\mathbf{d})\cdot\boldsymbol{\mu}_2
  \bigr),
  \label{eq:mu1} \\[0.35em]
  \boldsymbol{\mu}_2
  &=
  \boldsymbol{\alpha}_2
  \bigl(
    \mathbf{E}_0
    +
    \mathbf{M}(\mathbf{d})\cdot\boldsymbol{\mu}_1
  \bigr),
  \label{eq:mu2}
\end{align}
where $\mathbf{d}$ is the vector from body~2 to body~1 and $\mathbf{M}(\mathbf{d})$ is the
quasistatic dipole--dipole interaction tensor~\citep{jackson1999,gersten1980},
\begin{equation}
  \mathbf{M}(\mathbf{d})
  =
  \frac{3\,\hat{\mathbf{n}}\hat{\mathbf{n}}-\mathbf{I}}{d^3},
  \qquad
  \hat{\mathbf{n}}=\mathbf{d}/d,
  \label{eq:Mtensor}
\end{equation}
valid for $d\ll\lambda$ (no retardation).  Equation~\eqref{eq:mu1} states that the
molecular dipole feels $\mathbf{E}_0$ plus the field produced by $\boldsymbol{\mu}_2$;
Eq.~\eqref{eq:mu2} states the converse.

\paragraph{Step 2: eliminate $\boldsymbol{\mu}_2$ and define $\boldsymbol{\alpha}_{\mathrm{eff}}$.}
Substitute~\eqref{eq:mu2} into~\eqref{eq:mu1}:
\begin{equation}
  \boldsymbol{\mu}_1
  =
  \boldsymbol{\alpha}_1
  \Bigl[
    \mathbf{E}_0
    +
    \mathbf{M}\cdot\boldsymbol{\alpha}_2\cdot\mathbf{E}_0
    +
    \mathbf{M}\cdot\boldsymbol{\alpha}_2\cdot\mathbf{M}\cdot\boldsymbol{\mu}_1
  \Bigr].
\end{equation}
Collect terms that multiply $\boldsymbol{\mu}_1$:
\begin{equation}
  \boldsymbol{\mu}_1
  -
  \boldsymbol{\alpha}_1\cdot\mathbf{M}\cdot\boldsymbol{\alpha}_2\cdot\mathbf{M}\cdot\boldsymbol{\mu}_1
  =
  \boldsymbol{\alpha}_1\cdot\bigl(\mathbf{I}+\mathbf{M}\cdot\boldsymbol{\alpha}_2\bigr)\cdot\mathbf{E}_0,
\end{equation}
or
\begin{equation}
  \bigl(
    \mathbf{I}
    -
    \boldsymbol{\alpha}_1\cdot\mathbf{M}\cdot\boldsymbol{\alpha}_2\cdot\mathbf{M}
  \bigr)
  \boldsymbol{\mu}_1
  =
  \boldsymbol{\alpha}_1\cdot\bigl(\mathbf{I}+\mathbf{M}\cdot\boldsymbol{\alpha}_2\bigr)\cdot\mathbf{E}_0.
  \label{eq:mu1-lin}
\end{equation}
Inverting the operator on the left-hand side gives Gersten and Nitzan's effective
polarizability of the molecule~\citep{gersten1980},
\begin{equation}
  \boldsymbol{\mu}_1
  =
  \boldsymbol{\alpha}_{\mathrm{eff}}\cdot\mathbf{E}_0,
  \qquad
  \boldsymbol{\alpha}_{\mathrm{eff}}
  =
  \bigl(
    \mathbf{I}
    -
    \boldsymbol{\alpha}_1\cdot\mathbf{M}\cdot\boldsymbol{\alpha}_2\cdot\mathbf{M}
  \bigr)^{-1}
  \cdot
  \boldsymbol{\alpha}_1
  \cdot
  \bigl(
    \mathbf{I}
    +
    \mathbf{M}\cdot\boldsymbol{\alpha}_2
  \bigr).
  \label{eq:aeff-GN}
\end{equation}
Equation~\eqref{eq:aeff-GN} is their Eq.~(1.5): the molecule's response to the
\emph{incident} field already incorporates electromagnetic coupling to body~2.

\paragraph{Step 3: scalar reduction and identification of $G$ and $S$.}
Along a principal axis (radial or tangential), $\boldsymbol{\alpha}_1$, $\boldsymbol{\alpha}_2$,
and $\mathbf{M}$ reduce to scalars $\alpha_m$, $\alpha_2$, and $M$.  Then
\begin{equation}
  \alpha_{\mathrm{eff}}
  =
  \frac{\alpha_m\,(1+M\alpha_2)}{1-\alpha_m\,(M\alpha_2 M)}
  =
  \frac{\alpha_m\,(1+M\alpha_2)}{1-\alpha_m\,M^2\alpha_2}.
  \label{eq:aeff-scalar}
\end{equation}
Define
\begin{align}
  G
  &\equiv
  1 + M\alpha_2,
  \label{eq:G-def} \\[0.25em]
  S
  &\equiv
  M\,\alpha_2\,M
  =
  M^2\alpha_2.
  \label{eq:S-def}
\end{align}
Then~\eqref{eq:aeff-scalar} is exactly
\begin{equation}
  \alpha_{\mathrm{eff}}
  =
  \frac{\alpha_m\, G}{1 - \alpha_m\, S}.
  \label{eq:aeff-GS}
\end{equation}
The physical content of~\eqref{eq:G-def}--\eqref{eq:S-def} is transparent:
\begin{itemize}[leftmargin=1.4em]
  \item $G E_0$ is the field at the molecular site \emph{when the molecule is absent}:
        the incident field plus the field of body~2 driven by $\mathbf{E}_0$ alone
        (the local-field or ``lightning-rod'' contribution of Ref.~\citenum{gersten1980}).
  \item $S p$ is the additional field returned to the molecule when its own dipole
        $p$ polarizes body~2 (the image / self-interaction contribution).
\end{itemize}
Gersten and Nitzan emphasize the same decomposition: local-field enhancement
from polarizing the body in $\mathbf{E}_0$, and image enhancement from the body
responding to the molecular dipole~\citep{gersten1980}.

\paragraph{Step 4: total field at the molecule (equivalent one-site form).}
An equivalent and convenient statement is obtained by writing the total field
acting on the molecule as
\begin{equation}
  E_{\mathrm{tot}}
  =
  G\, E_{\mathrm{inc}}
  +
  S\, p,
  \qquad
  p
  =
  \alpha_m\, E_{\mathrm{tot}}.
  \label{eq:Etot}
\end{equation}
(The first term is present even if $p=0$; the second is linear in $p$.)
Eliminating $E_{\mathrm{tot}}$,
\begin{equation}
  p
  =
  \alpha_m\bigl(G E_{\mathrm{inc}} + S p\bigr)
  =
  \alpha_m G\, E_{\mathrm{inc}}
  +
  \alpha_m S\, p,
\end{equation}
so
\begin{equation}
  p - \alpha_m S\, p
  =
  \alpha_m G\, E_{\mathrm{inc}}
  \qquad\Rightarrow\qquad
  p
  =
  \frac{\alpha_m G}{1-\alpha_m S}\, E_{\mathrm{inc}},
\end{equation}
which is~\eqref{eq:alphaeff}.  Setting $S=0$ recovers the pure local-field limit
$\alpha_{\mathrm{eff}}=\alpha_m G$ (image neglected), as discussed in
Ref.~\citenum{gersten1980} when the inverse matrices associated with the image
are dropped.

\paragraph{Step 5: dipole tensor eigenvalues for parallel and perpendicular.}
For a unit dipole oriented along $\hat{\mathbf{n}}=\hat{\mathbf{r}}$ (radial /
parallel to the sphere--molecule axis),
\begin{equation}
  \mathbf{M}\cdot\hat{\mathbf{r}}
  =
  \frac{2}{d^3}\,\hat{\mathbf{r}}
  \qquad\Rightarrow\qquad
  M_\parallel = \frac{2}{d^3}.
\end{equation}
For a unit dipole tangential to the sphere (perpendicular orientation),
\begin{equation}
  \mathbf{M}\cdot\hat{\boldsymbol{\theta}}
  =
  -\frac{1}{d^3}\,\hat{\boldsymbol{\theta}}
  \qquad\Rightarrow\qquad
  M_\perp = -\frac{1}{d^3}.
\end{equation}
If body~2 is replaced by a sphere whose \emph{dipole} polarizability volume is
$\beta_1$ from Eq.~\eqref{eq:beta1}, then~\eqref{eq:G-def}--\eqref{eq:S-def} give
\begin{align}
  G_\parallel
  &=
  1 + M_\parallel\beta_1
  =
  1 + 2\,\frac{\beta_1}{d^3},
  &
  S_\parallel^{(\ell=1)}
  &=
  M_\parallel^2\beta_1
  =
  4\,\frac{\beta_1}{d^6},
  \label{eq:GSdip-par} \\[0.35em]
  G_\perp
  &=
  1 + M_\perp\beta_1
  =
  1 - \frac{\beta_1}{d^3},
  &
  S_\perp^{(\ell=1)}
  &=
  M_\perp^2\beta_1
  =
  \frac{\beta_1}{d^6},
  \label{eq:GSdip-perp}
\end{align}
which matches Eqs.~\eqref{eq:Gpar}--\eqref{eq:Gperp} for $G$ and the $\ell=1$
truncation of the multipole series~\eqref{eq:Spar}--\eqref{eq:Sperp} for $S$
[since $(\ell+1)^2\beta_1/d^{2\ell+4}$ at $\ell=1$ is $4\beta_1/d^6$].

\paragraph{Step 6: multipolar sphere (beyond a single point polarizability).}
A real sphere is not a pure point dipole: a molecular source at finite $d$
excites all multipoles $\ell\ge 1$.  The field of the sphere at the molecular
site remains linear in the molecular dipole, so one may still write
$E_{\mathrm{refl}}=S p$, but $S$ is no longer $M^2\beta_1$.  Matching boundary
conditions for a point dipole outside a dielectric sphere (or spheroid) yields
the multipole series of Sec.~\ref{sec:fields},
Eqs.~\eqref{eq:Spar}--\eqref{eq:Sperp}~\citep{gersten1981,ruppin1982,novotny2012}.
Likewise, the field of the sphere driven by $\mathbf{E}_{\mathrm{inc}}$ alone is
still of the form $G E_{\mathrm{inc}}$ with $G$ given by the dipole channel
\eqref{eq:Gpar}--\eqref{eq:Gperp} for a uniform incident field.

In Gersten and Nitzan's 1981 treatment of a molecule outside a small dielectric
particle, the same physics appears as an \emph{image enhancement factor}
$A(\parallel)$, $A(\perp)$ built from multipole sums; molecular amplitudes
acquire factors $(1-A)^{-1}$~\citep{gersten1981}.  That factor is the same
renormalization as $(1-\alpha_m S)^{-1}$ once $A$ is identified with
$\alpha_m S$ in the present notation.  Thus:
\begin{itemize}[leftmargin=1.4em]
  \item the \emph{structure} $\alpha_{\mathrm{eff}}=\alpha_m G/(1-\alpha_m S)$
        follows from the two-body classical electrodynamics of
        Ref.~\citenum{gersten1980};
  \item the \emph{explicit multipole forms} of $G$ and $S$ for a sphere (or
        spheroid) are those of the quasistatic boundary-value problem developed
        for molecule--particle systems in Ref.~\citenum{gersten1981}.
\end{itemize}

\paragraph{Step 7: summary of the working formula.}
Collecting the pieces used in this model:
\begin{enumerate}[leftmargin=1.4em]
  \item Choose orientation $\parallel$ or $\perp$ and host permittivity
        $\varepsilon_b$.
  \item Compute $\beta_\ell(\omega)$ from~\eqref{eq:beta} and form
        $G(\omega)$ and $S(\omega)$ from~\eqref{eq:Gpar}--\eqref{eq:Sperp}.
  \item Insert the bare molecular polarizability $\alpha_m(\omega)$
        from~\eqref{eq:alpham}.
  \item Evaluate
        \[
          \alpha_{\mathrm{eff}}(\omega)
          =
          \frac{\alpha_m(\omega)\, G(\omega)}{1-\alpha_m(\omega)\, S(\omega)}.
        \]
\end{enumerate}
The numerator enhances the drive of the molecule by the nanoparticle local field;
the denominator shifts and reshapes the molecular resonance through image
(self-)interaction with the sphere.  Both ingredients are required for a consistent
classical treatment of a polarizable molecule outside a polarizable particle.

% =============================================================================
\section{Absorption spectrum}
\label{sec:abs}
% =============================================================================

In linear response, dissipation associated with the induced dipole driven by a
real monochromatic field is proportional to the imaginary part of the
polarizability~\citep{jackson1999,novotny2012}.  An absorption-like lineshape consistent with the Fourier
transform of a causal dipole response is
\begin{equation}
  A(\omega)
  =
  -\omega\,
  \operatorname{Im}
  \alpha_{\mathrm{eff}}(\omega).
  \label{eq:A}
\end{equation}
Overall scale is fixed by $\alpha_0$ and is often removed by peak normalization,
\begin{equation}
  A_{\mathrm{norm}}(\omega)
  =
  \frac{A(\omega)}{\max_{\omega}|A(\omega)|},
\end{equation}
optionally after a global sign flip if the dominant extremum is negative (a
convention that depends on Fourier phase conventions when $A$ is reconstructed
from time-domain data).\\

\noindent
Important interpretive points:
\begin{itemize}[leftmargin=1.4em]
  \item $A(\omega)$ is the \emph{molecular-channel} spectrum relative to
        $E_{\mathrm{inc}}$.  It is not the extinction of the composite particle.
  \item Plasmonic structure enters only through $G(\omega)$ and $S(\omega)$, and
        therefore through $\varepsilon(\omega)$.
  \item A second peak near $\omega_{\mathrm{LSPR}}$ appears only if
        $\operatorname{Im}(\alpha_m G)$ or the hybrid poles of
        $1-\alpha_m S=0$ place appreciable weight there---not because nanoparticle
        absorption is plotted directly.
\end{itemize}

% =============================================================================
\section{Parameterization of gold: Meep material \texttt{Au}}
\label{sec:au}
% =============================================================================

\subsection{Source and functional form}

The dielectric function of the sphere is taken from the bulk gold model
distributed as the material \texttt{Au} in the Meep electromagnetic library
(\texttt{meep.materials.Au})~\citep{oskooi2010,meepdocs}.  That implementation follows the analytic
Drude--Lorentz fit of Raki\'{c} \emph{et al.}~\citep{rakic1998} to tabulated optical constants of
gold over a broad infrared--ultraviolet window.\\

\noindent
In frequency units where susceptibilities are written with energies $E=\hbar\omega$
in electronvolts, Meep's Lorentzian and Drude poles~\citep{oskooi2010} reduce to
\begin{align}
  \chi_{\mathrm{L}}(E)
  &=
  \frac{\mathcal{A}}{E_0^2 - E^2 - i\Gamma E},
  \label{eq:chiL} \\[0.35em]
  \chi_{\mathrm{D}}(E)
  &=
  \frac{\mathcal{A}}{-E^2 - i\Gamma E},
  \label{eq:chiD}
\end{align}
with amplitude $\mathcal{A}=\sigma E_0^2$ related to Meep's dimensionless
oscillator strength $\sigma$ and pole energy $E_0$.  The full permittivity is
\begin{equation}
  \varepsilon(E)
  =
  \varepsilon_\infty
  +
  \chi_{\mathrm{D}}(E)
  +
  \sum_{j=1}^{N_{\mathrm{L}}}
  \chi_{\mathrm{L},j}(E).
  \label{eq:eps-au}
\end{equation}

\subsection{Raki\'{c} parameters for \texttt{Au}}

For \texttt{Au} one has $\varepsilon_\infty=1$ and a plasma energy
$\hbar\omega_p=9.03\,\mathrm{eV}$.  The Drude weight is set by the free-electron
oscillator strength $f_0=0.760$ and damping $\Gamma_0=0.053\,\mathrm{eV}$,
\begin{equation}
  \mathcal{A}_{\mathrm{D}}
  =
  f_0\,(\hbar\omega_p)^2
  =
  0.760\times(9.03\,\mathrm{eV})^2,
\end{equation}
as tabulated by Raki\'{c} \emph{et al.}~\citep{rakic1998}.
Five Lorentzian interband (and multi-electron) terms complete the fit.  In the
Raki\'{c} parameterization each Lorentzian carries strength $f_j$ at energy
$E_j$ with damping $\Gamma_j$, and
\begin{equation}
  \mathcal{A}_j = f_j\,(\hbar\omega_p)^2,
\end{equation}
which is equivalent to Meep's $\sigma_j=f_j\omega_p^2/E_j^2$ after conversion of
units.
\begin{center}
\begin{tabular}{@{}clccc@{}}
\toprule
$j$ & Type & $E_j$ (eV) & $\Gamma_j$ (eV) & $f_j$ \\
\midrule
0 & Drude   & ---   & $0.053$ & $0.760$ \\
1 & Lorentz & $0.415$ & $0.241$ & $0.024$ \\
2 & Lorentz & $0.830$ & $0.345$ & $0.010$ \\
3 & Lorentz & $2.969$ & $0.870$ & $0.071$ \\
4 & Lorentz & $4.304$ & $2.494$ & $0.601$ \\
5 & Lorentz & $13.32$ & $2.214$ & $4.384$ \\
\bottomrule
\end{tabular}
\end{center}

\noindent
These values are those encoded in Meep's \texttt{Au} material following
Ref.~\citenum{rakic1998}.  The stated validity range of the fit spans roughly
$0.25$--$6.2\,\mu\mathrm{m}$~\citep{rakic1998}.

\subsection{Consequences for the sphere resonance}

With this $\varepsilon(E)$, the quasistatic Fr\"{o}hlich energy defined by
$\operatorname{Re}\varepsilon(E)=-2\varepsilon_b$ depends on the host~\citep{bohren1983,maier2007}.  In vacuum
($\varepsilon_b=1$) one finds
\begin{equation}
  E_{\mathrm{LSPR}}^{\mathrm{(qs,vac)}}
  \approx
  2.60\,\mathrm{eV},
\end{equation}
whereas in water ($n=1.33$, $\varepsilon_b\approx 1.769$, target
$\operatorname{Re}\varepsilon\approx-3.54$)
\begin{equation}
  E_{\mathrm{LSPR}}^{\mathrm{(qs,water)}}
  \approx
  2.42\,\mathrm{eV}.
\end{equation}
Finite-size retardation and radiative damping further shift the extinction peak
of a real nanoparticle relative to these bulk quasistatic estimates~\citep{bohren1983,maier2007}.  Within the
present model, plasmonic structure in $G$ and $S$ is locked to
Eqs.~\eqref{eq:eps-au}--\eqref{eq:frohlich} with the table above and the chosen
$\varepsilon_b$.

\subsection{Unit conversion from Meep frequencies}

Meep stores pole frequencies $f$ in units of $c/(1\,\mu\mathrm{m})$ when the
characteristic length scale is one micrometre~\citep{oskooi2010,meepdocs}.  Conversion to electronvolts uses
\begin{equation}
  E[\mathrm{eV}]
  =
  \frac{f}{\eta},
  \qquad
  \eta
  =
  \frac{1}{1.23984193}\,\mu\mathrm{m}^{-1}\,\mathrm{eV}^{-1}.
\end{equation}
After this conversion, Eqs.~\eqref{eq:chiL}--\eqref{eq:chiD} are free of
explicit $2\pi$ factors because Meep's angular-frequency convention cancels
between numerator and denominator of each susceptibility.

% =============================================================================
\section{Parameterization of sodium}
\label{sec:na}
% =============================================================================

\subsection{Physical transition}

Atomic sodium is treated as a single effective optical transition corresponding
to the $D$ lines ($3s\to 3p$), represented by the Lorentz oscillator
\eqref{eq:alpham}.  The sodium $D_1$ and $D_2$ resonances lie near
$589.6\,\mathrm{nm}$ and $589.0\,\mathrm{nm}$~\citep{nistasd}, often combined into one
effective optical peak in reduced models.  In the model the three scalar parameters are:
\begin{itemize}[leftmargin=1.4em]
  \item $\omega_m$ --- resonance energy of the absorption peak of the isolated
        atom (or of the chosen electronic-structure calculation for Na);
  \item $\gamma_m$ --- phenomenological damping (lifetime, dephasing, and any
        artificial broadening);
  \item $\alpha_0$ --- static polarizability volume fixing the integrated
        oscillator strength of the classical dipole.
\end{itemize}

\subsection{Typical values}

A representative optical resonance energy for Na in the visible is of order
\begin{equation}
  \omega_m \sim 2.1\,\mathrm{eV},
\end{equation}
consistent with the sodium $D$ lines near $589\,\mathrm{nm}$
($2.104\,\mathrm{eV}$) when a single effective peak is retained~\citep{nistasd}.  The natural
radiative linewidth is far smaller than typical condensed-phase or numerical
broadenings; $\gamma_m$ is therefore an effective parameter (often
$10^{-2}$--$10^{-1}\,\mathrm{eV}$ in reduced models).\\

\noindent
The static electric dipole polarizability of atomic sodium is known accurately
from theory and experiment~\citep{ekstrom1995,mitroy2010},
\begin{equation}
  \alpha_{\mathrm{Na}}(0)
  \approx
  24.1\,\text{\AA}^3
  =
  0.0241\,\mathrm{nm}^3.
\end{equation}
In Eq.~\eqref{eq:alpham} one may set $\alpha_0$ to this value if the Lorentzian
is understood as carrying the full static polarizability, or treat $\alpha_0$ as
an effective resonant weight when only the optical $D$-line contribution is of
interest.  Increasing $\alpha_0$ strengthens both the bare absorption and the
back-action factor $\alpha_m S$ near the sphere.

\subsection{Coupling geometry}

With Na represented as a point dipole, the only geometric parameters entering
the hybrid response are $(a,d)$ and the orientation.  No internal structure of
the atom (orbital shape, multipolar transition density) is retained beyond
$\alpha_m(\omega)$.  The model is therefore appropriate when the surface gap is
large compared with the atomic size, and becomes increasingly schematic for
sub-nanometre gaps where nonlocal response of the metal and finite orbital
extent matter~\citep{ciraci2012,zhu2016}.

% =============================================================================
\section{Summary of the closed-form model}
\label{sec:summary}
% =============================================================================

Given $\varepsilon(\omega)$ for gold (Sec.~\ref{sec:au}), host index
$n$ (hence $\varepsilon_b=n^2$; default $n=1.33$), molecular parameters
$(\omega_m,\gamma_m,\alpha_0)$ for sodium (Sec.~\ref{sec:na}), geometry
$(a,d)$, and orientation, the absorption spectrum is obtained as follows:
\begin{enumerate}[leftmargin=1.4em]
  \item Evaluate $\varepsilon(\omega)$ from the Raki\'{c}/\texttt{Au}
        Drude--Lorentz expansion~\citep{rakic1998,oskooi2010}.
  \item Form $\beta_\ell(\omega)$ from Eq.~\eqref{eq:beta} with the host
        $\varepsilon_b$~\citep{bohren1983}.
  \item Compute $G(\omega)$ and $S(\omega)$ from
        Eqs.~\eqref{eq:Gpar}--\eqref{eq:Sperp}~\citep{gersten1980,gersten1981}.
  \item Form $\alpha_m(\omega)$ from Eq.~\eqref{eq:alpham}.
  \item Construct $\alpha_{\mathrm{eff}}(\omega)$ from
        Eq.~\eqref{eq:alphaeff}~\citep{gersten1980}.
  \item Evaluate $A(\omega)=-\omega\,\operatorname{Im}\alpha_{\mathrm{eff}}(\omega)$.
\end{enumerate}

\noindent
All steps are algebraic once $\varepsilon(\omega)$ is known; the only numerical
tasks are evaluation on a frequency grid and truncation of the multipole sum at
convergent $\ell_{\max}$.

% =============================================================================
\section{Limitations}
\label{sec:limits}
% =============================================================================

\begin{itemize}[leftmargin=1.4em]
  \item \textbf{Quasistatic sphere.}
        Retardation, radiative damping, and higher Mie coefficients of the
        driven sphere are omitted~\citep{bohren1983,maier2007}.  Extinction peaks of finite spheres therefore
        need not coincide with the Fr\"{o}hlich energy
        $\operatorname{Re}\varepsilon=-2\varepsilon_b$.
  \item \textbf{Non-dispersive host.}
        Water is treated as a real constant $n=1.33$.  Dispersion and absorption
        of the solvent are neglected.
  \item \textbf{Local bulk dielectric.}
        Electron spill-out, nonlocality, and surface scattering in the metal are
        absent from Raki\'{c} bulk $\varepsilon(\omega)$~\citep{rakic1998,ciraci2012}.
  \item \textbf{Point molecule.}
        Finite orbital size and multipolar transitions of Na are neglected.
  \item \textbf{Linear classical response.}
        No saturation, quantum strong-coupling corrections beyond classical
        hybridization through $S$, or non-Markovian metal continuum beyond the
        fitted poles~\citep{torma2015}.
  \item \textbf{Single Na oscillator.}
        Additional electronic transitions of sodium or a molecule are omitted
        unless extra Lorentzians are added to $\alpha_m$.
\end{itemize}

% =============================================================================
\begin{thebibliography}{99}

\bibitem{mie1908}
G.~Mie,
``Beitr\"{a}ge zur Optik tr\"{u}ber Medien, speziell kolloidaler Metall\"{o}sungen,''
\emph{Ann.\ Phys.} \textbf{330}, 377--445 (1908).
\doiref{10.1002/andp.19083300302}.

\bibitem{frohlich1958}
H.~Fr\"{o}hlich,
\emph{Theory of Dielectrics},
2nd~ed.\ (Oxford University Press, Oxford, 1958).
% No DOI for this edition; catalog record often listed under ISBN~978-0-19-851379-7.

\bibitem{bohren1983}
C.~F.~Bohren and D.~R.~Huffman,
\emph{Absorption and Scattering of Light by Small Particles}
(Wiley, New York, 1983).
\doiref{10.1002/9783527618156}.

\bibitem{gersten1980}
J.~Gersten and A.~Nitzan,
``Electromagnetic theory of enhanced Raman scattering by molecules adsorbed on
rough surfaces,''
\emph{J.\ Chem.\ Phys.} \textbf{73}, 3023--3037 (1980).
\doiref{10.1063/1.440560}.

\bibitem{gersten1981}
J.~Gersten and A.~Nitzan,
``Spectroscopic properties of molecules interacting with small dielectric
particles,''
\emph{J.\ Chem.\ Phys.} \textbf{75}, 1139--1152 (1981).
\doiref{10.1063/1.442161}.

\bibitem{ruppin1982}
R.~Ruppin,
``Optical properties of small metal spheres,''
\emph{Phys.\ Rev.\ B} \textbf{26}, 3440--3444 (1982).
\doiref{10.1103/PhysRevB.26.3440}.

\bibitem{jackson1999}
J.~D.~Jackson,
\emph{Classical Electrodynamics},
3rd~ed.\ (Wiley, New York, 1999).
% No DOI for this edition; ISBN~978-0-471-30932-1.

\bibitem{maier2007}
S.~A.~Maier,
\emph{Plasmonics: Fundamentals and Applications}
(Springer, New York, 2007).
\doiref{10.1007/0-387-37825-1}.

\bibitem{novotny2012}
L.~Novotny and B.~Hecht,
\emph{Principles of Nano-Optics},
2nd~ed.\ (Cambridge University Press, Cambridge, 2012).
\doiref{10.1017/CBO9780511794193}.

\bibitem{rakic1998}
A.~D.~Raki\'{c}, A.~B.~Djuri\v{s}i\'{c}, J.~M.~Elazar, and M.~L.~Majer,
``Optical properties of metallic films for vertical-cavity optoelectronic
devices,''
\emph{Appl.\ Opt.} \textbf{37}, 5271--5283 (1998).
\doiref{10.1364/AO.37.005271}.

\bibitem{oskooi2010}
A.~F.~Oskooi, D.~Roundy, M.~Ibanescu, P.~Bermel, J.~D.~Joannopoulos, and
S.~G.~Johnson,
``{MEEP}: A flexible free-software package for electromagnetic simulations by
the {FDTD} method,''
\emph{Comput.\ Phys.\ Commun.} \textbf{181}, 687--702 (2010).
\doiref{10.1016/j.cpc.2009.11.008}.

\bibitem{meepdocs}
A.~F.~Oskooi, D.~Roundy, M.~Ibanescu, P.~Bermel, J.~D.~Joannopoulos, and
S.~G.~Johnson,
``{MEEP}: A flexible free-software package for electromagnetic simulations by
the {FDTD} method,''
\emph{Comput.\ Phys.\ Commun.} \textbf{181}, 687--702 (2010);
materials library documentation at \url{https://meep.readthedocs.io/}.
\doiref{10.1016/j.cpc.2009.11.008}.

\bibitem{nistasd}
A.~Kramida, Yu.~Ralchenko, J.~Reader, and {NIST ASD Team},
\emph{{NIST} Atomic Spectra Database} (ver.\ 5.x),
National Institute of Standards and Technology, Gaithersburg, MD.
\doiref{10.18434/T4W30F}.

\bibitem{ekstrom1995}
C.~R.~Ekstrom, J.~Schmiedmayer, M.~S.~Chapman, T.~D.~Hammond, and D.~E.~Pritchard,
``Measurement of the electric polarizability of sodium with an atom interferometer,''
\emph{Phys.\ Rev.\ A} \textbf{51}, 3883--3888 (1995).
\doiref{10.1103/PhysRevA.51.3883}.

\bibitem{mitroy2010}
J.~Mitroy, M.~S.~Safronova, and C.~W.~Clark,
``Theory and applications of atomic and ionic polarizabilities,''
\emph{J.\ Phys.\ B: At.\ Mol.\ Opt.\ Phys.} \textbf{43}, 202001 (2010).
\doiref{10.1088/0953-4075/43/20/202001}.

\bibitem{ciraci2012}
C.~Cirac\`{i}, R.~T.~Hill, J.~J.~Mock, Y.~Urzhumov, A.~I.~Fern\'{a}ndez-Dom\'{i}nguez,
S.~A.~Maier, J.~B.~Pendry, A.~Chilkoti, and D.~R.~Smith,
``Probing the ultimate limits of plasmonic enhancement,''
\emph{Science} \textbf{337}, 1072--1074 (2012).
\doiref{10.1126/science.1224823}.

\bibitem{zhu2016}
W.~Zhu, R.~Esteban, A.~G.~Borisov, J.~J.~Baumberg, P.~Nordlander, H.~J.~Lezec,
J.~Aizpurua, and K.~B.~Crozier,
``Quantum mechanical effects in plasmonic structures with subnanometre gaps,''
\emph{Nat.\ Commun.} \textbf{7}, 11495 (2016).
\doiref{10.1038/ncomms11495}.

\bibitem{torma2015}
P.~T\"{o}rm\"{a} and W.~L.~Barnes,
``Strong coupling between surface plasmon polaritons and emitters: a review,''
\emph{Rep.\ Prog.\ Phys.} \textbf{78}, 013901 (2015).
\doiref{10.1088/0034-4885/78/1/013901}.

\end{thebibliography}

\end{document}
